% =====================================================================
% 7. CONCLUSIONS
% =====================================================================
\section{Conclusions}

\subsection{Conclusions related to the study}

\begin{itemize}
  \item \textbf{H1 (classification) is supported.} Age and 20 point mutations separate LGG from GBM
        with out-of-fold AUCs between \ClfMinAUC{} and \ClfMaxAUC{} under nested cross-validation. Penalised logistic
        regression and XGBoost are statistically indistinguishable from the best random forest;
        naive Bayes, the SVM and the pruned tree are slightly worse in paired comparisons.
  \item \textbf{H2 (regression) is supported with a caveat.} Without the grade, the profile
        explains a share of the variance of age between \RegMinRTwo{} and \RegMaxRTwo{} ($R^2$), with all
        models equivalent; but the predictions
        mostly recover the grade through IDH1, IDH2 and ATRX, so the mutations explain age largely
        through the grade.
  \item \textbf{H3 (unsupervised) is supported.} Without labels, $k$-means and Ward's method
        separate an IDH-wildtype and an IDH-mutant group, related to but not identical with the
        grade; a secondary three-cluster solution separates two IDH-mutant profiles. Jaccard-based
        hierarchical clustering fails on these data because of patients without any mutation.
  \item \textbf{The labels limit the study.} The grades come from pre-2021 criteria and several
        markers of the current classification are missing from the data; the cohort is mostly white
        and only internally validated. The models describe this cohort, not current clinical
        practice.
\end{itemize}

\subsection{Conclusions on what you have learned}

The most useful lessons were methodological. First, a large part of good modelling happens before
the models: the histological diagnosis would have produced a perfect but meaningless classifier, and
only checking how the target was defined revealed it. Second, the choice of a distance or a linkage
is a modelling decision with consequences: the Jaccard distance is the natural choice for sparse
binary data, but it made hierarchical clustering degenerate. Third, ``which model is best'' often has
no answer at the available sample size, and saying so with intervals is more honest than ranking
point estimates; how the comparison is made (marginal intervals or paired differences) must be fixed
in advance. Fourth, a model can rank patients well while giving unreliable probabilities or
coefficients (naive Bayes, logistic regression under quasi-separation), so evaluation must match the
intended use. Finally, the AI assistants used for literature search produced wrong or misattributed
sources that were only caught by cross-checking (Appendix B).
